\documentclass[11pt]{article}
\usepackage{fontspec,amsmath,amssymb,amsthm,hyperref}
\usepackage[margin=1in]{geometry}
\setlength{\emergencystretch}{3em}
\title{A weak forward-equation interface and a conditional stationarity route}
\author{LangevinLean mathematical exposition}
\date{}
\begin{document}\maketitle
\noindent Source module: \texttt{UnadjustedLangevin/FokkerPlanck.lean}. Compile this proof document with XeLaTeX or LuaLaTeX. The accompanying HTML places the complete Lean declarations beside every proof.


The weak Fokker–Planck equation is formulated by integrating \(\mathcal L\varphi\) against time-dependent probability measures. The final theorem in this module assumes uniqueness of this weak equation explicitly; the main ULA development proves stationarity by the separate Metropolis approximation route.


\section{Specify the weak solution class}


A path \(\rho_t\) is a weak solution when each \(\rho_t\) is a probability measure with finite second moment, compactly supported \(C^2\) tests are integrable and have continuous averages, generator averages are integrable in space and time, and

\[\int\varphi\,d\rho_t=\int\varphi\,d\rho_0+\int_0^t\int\mathcal L\varphi\,d\rho_s\,ds.\]

The constant path \(\rho_t=\pi\) satisfies every field using Gibbs normalization, its second moment, compact-support integrability and \(\int\mathcal L\varphi\,d\pi=0\).



\noindent Lean declarations: \texttt{IsWeakFokkerPlanckSolution}, \texttt{constant\_gibbs\_isWeakFokkerPlanckSolution}.

\section{Use uniqueness only as an explicit premise}


Define weak uniqueness to mean that any two paths in this solution class with the same initial law agree on \([0,T]\). Apply this property to the supplied solution \(\rho\) and the already proved constant Gibbs solution. Their equality at time zero yields \(\rho_t=\pi\). No theorem here derives the uniqueness premise from smoothness.



\noindent Lean declarations: \texttt{WeakFokkerPlanckUnique}, \texttt{gibbs\_stationary\_of\_weakFokkerPlanckUniqueness}.

\end{document}
